\subsection*{C++ 主机：形状返回}

主机侧收到的结果拥有自己的缓冲区，具有值语义，便于在调用结束后继续使用。

\begin{lstlisting}[language=croqtile,caption={返回带形状的缓冲区}]
__co__ auto foo() {
  f32 [7, 16] output;
  // ...
  return output;  // semantically copy
}
void entry() {
  auto result = foo(); // move the 'spanned_data' to caller
  // ...
}
\end{lstlisting}

示例中，\co@output@ 作为结果返回，主机侧类型为 \co@croq::spanned_data@。实现通过返回值优化避免复制整块缓冲区。简化定义如下。

\begin{lstlisting}[language=croqtile,caption={主机侧结果包装器}]
template <typename T, size_t Rank>
class spanned_data {
  std::unique_ptr<T[]> ptr; // move only pointers
  size_t dims[Rank];        // ranked dimensions
};
\end{lstlisting}

\co@spanned_view@ 只引用数据，\co@spanned_data@ 则管理缓冲区的所有权。返回值优化允许在交付结果时省去数据复制。

主机侧可以用 \co@.rank()@、\co@.shape()@ 查询结果形状，而不必依赖核函数内部的声明。
